\documentclass[10pt,a4paper]{amsart}
\usepackage{amsmath,amssymb,amsthm,graphicx}
\usepackage{verbatim}
\newenvironment{reference}{\comment}{\endcomment}
%\newenvironment{reference}{}{}
\newenvironment{slogan}{\comment}{\endcomment}
\newenvironment{history}{\comment}{\endcomment}

% For commutative diagrams we use Xy-pic
\usepackage[all]{xy}

% We use 2cell for 2-commutative diagrams.
\xyoption{2cell}
% Theorem environments.
%
\theoremstyle{plain}
\newtheorem{theorem}[subsection]{Theorem}
\newtheorem{proposition}[subsection]{Proposition}
\newtheorem{lemma}[subsection]{Lemma}

\theoremstyle{definition}
\newtheorem{definition}[subsection]{Definition}
\newtheorem{example}[subsection]{Example}
\newtheorem{exercise}[subsection]{Exercise}
\newtheorem{situation}[subsection]{Situation}

\theoremstyle{remark}
\newtheorem{remark}[subsection]{Remark}
\newtheorem{remarks}[subsection]{Remarks}

\numberwithin{equation}{subsection}

% Macros
%
\def\lim{\mathop{\mathrm{lim}}\nolimits}
\def\colim{\mathop{\mathrm{colim}}\nolimits}
\def\Spec{\mathop{\mathrm{Spec}}}
\def\Hom{\mathop{\mathrm{Hom}}\nolimits}
\def\Ext{\mathop{\mathrm{Ext}}\nolimits}
\def\SheafHom{\mathop{\mathcal{H}\!\mathit{om}}\nolimits}
\def\SheafExt{\mathop{\mathcal{E}\!\mathit{xt}}\nolimits}
\def\Sch{\mathit{Sch}}
\def\Mor{\mathop{\mathrm{Mor}}\nolimits}
\def\Ob{\mathop{\mathrm{Ob}}\nolimits}
\def\Sh{\mathop{\mathit{Sh}}\nolimits}
\def\NL{\mathop{N\!L}\nolimits}
\def\CH{\mathop{\mathrm{CH}}\nolimits}
\def\proetale{{pro\text{-}\acute{e}tale}}
\def\etale{{\acute{e}tale}}
\def\QCoh{\mathit{QCoh}}
\def\Ker{\mathop{\mathrm{Ker}}}
\def\Im{\mathop{\mathrm{Im}}}
\def\Coker{\mathop{\mathrm{Coker}}}
\def\Coim{\mathop{\mathrm{Coim}}}

% Boxtimes
%
\DeclareMathSymbol{\boxtimes}{\mathbin}{AMSa}{"02}

%
% Macros for moduli stacks/spaces
%
\def\QCohstack{\mathcal{QC}\!\mathit{oh}}
\def\Cohstack{\mathcal{C}\!\mathit{oh}}
\def\Spacesstack{\mathcal{S}\!\mathit{paces}}
\def\Quotfunctor{\mathrm{Quot}}
\def\Hilbfunctor{\mathrm{Hilb}}
\def\Curvesstack{\mathcal{C}\!\mathit{urves}}
\def\Polarizedstack{\mathcal{P}\!\mathit{olarized}}
\def\Complexesstack{\mathcal{C}\!\mathit{omplexes}}
% \Pic is the operator that assigns to X its picard group, usage \Pic(X)
% \Picardstack_{X/B} denotes the Picard stack of X over B
% \Picardfunctor_{X/B} denotes the Picard functor of X over B
\def\Pic{\mathop{\mathrm{Pic}}\nolimits}
\def\Picardstack{\mathcal{P}\!\mathit{ic}}
\def\Picardfunctor{\mathrm{Pic}}
\def\Deformationcategory{\mathcal{D}\!\mathit{ef}}
\usepackage{cleveref}
\crefname{originallemma}{original lemma}{original lemmas}
\Crefname{originallemma}{Original lemma}{Original lemmas}
\makeatletter
\newlabel{etale-cohomology-lemma-directed-colimit-cohomology}{{03Q6}{source}{}{sourcepointer-1}{}}
\newlabel{etale-cohomology-lemma-directed-colimit-cohomology@cref}{{[originallemma][0][]03Q6}{[1][0][]source}{}{}{}}
\makeatother
\crefname{originalreference}{original reference}{original references}
\Crefname{originalreference}{Original reference}{Original references}
\makeatletter
\newlabel{lemma-describe-pullback-from-etale}{{0GLZ}{source}{}{sourcepointer-2}{}}
\newlabel{lemma-describe-pullback-from-etale@cref}{{[originalreference][0][]0GLZ}{[1][0][]source}{}{}{}}
\makeatother
\crefname{originalreference}{original reference}{original references}
\Crefname{originalreference}{Original reference}{Original references}
\makeatletter
\newlabel{lemma-limit-pullback}{{099S}{source}{}{sourcepointer-3}{}}
\newlabel{lemma-limit-pullback@cref}{{[originalreference][0][]099S}{[1][0][]source}{}{}{}}
\makeatother
\crefname{originalreference}{original reference}{original references}
\Crefname{originalreference}{Original reference}{Original references}
\makeatletter
\newlabel{proposition-weakly-etale}{{097Z}{source}{}{sourcepointer-4}{}}
\newlabel{proposition-weakly-etale@cref}{{[originalreference][0][]097Z}{[1][0][]source}{}{}{}}
\makeatother
\crefname{originalsection}{original section}{original sections}
\Crefname{originalsection}{Original section}{Original sections}
\makeatletter
\newlabel{sites-section-pullback}{{06UM}{source}{}{sourcepointer-5}{}}
\newlabel{sites-section-pullback@cref}{{[originalsection][0][]06UM}{[1][0][]source}{}{}{}}
\makeatother
\title{The pro-etale pushforward of the pullback of an etale sheaf is the pullback of the etale pushforward}
\author{The Stacks Project Authors}
\date{stacks-project:proetale.tex:local-snapshot-20261012; GNU FDL 1.2 or later}
\begin{document}
\maketitle
\noindent\textit{Editorial note.} All delivered mathematical text is the author's own, copied from the pinned Stacks Project chapter file proetale.tex; no statement or proof step was rewritten, reordered or removed.\par\medskip
\noindent\textit{Editorial note.} The counted proof is the author's complete argument: the reduction to an affine base and to limits of affine etale schemes, the identification of the two sides of the comparison map through the description of the pullback, and the closing computation of the colimit of the etale cohomology groups.\par\medskip
\noindent\textit{Editorial note.} The label lemma-morphism-comparison names the counted statement, and the proof environment immediately following it in the source is delivered verbatim as its proof.\par\medskip
\noindent\textit{Editorial note.} The counted statement is carried by the author's own lemma environment in the chapter source.\par\medskip
\noindent\textit{Editorial note.} Pre-open gate: python operations/agent-authoring/tools/precard\_gate.py --source proetale.tex --label lemma-morphism-comparison --start 4912 --end 4965 --local-id 3583 returned RESULT=OK (L3583 free, occupancy FREE, public interval clear). Candidate taken from sources/stacks-proetale/candidates-halfA.json.\par\medskip
\noindent\textit{Editorial note.} The statement and proof are copied from the pinned chapter file; the cited inputs are kept as external pointers with their canonical Stacks tags.\par\medskip
\noindent\textit{Allowed standard prerequisites (editorial).}\begin{itemize}
\item Description of the pullback along the pro-etale to etale comparison. Same chapter, the lemma describing the pullback of an etale sheaf to the pro-etale site as the pullback along the projection. Supplies the identification of the sections of the pullback over X x\_Y V with sections of the pullback of F along the projection to X.
\item Cohomology of a directed colimit of schemes. Etale Cohomology, the lemma computing the etale cohomology of a scheme written as a directed limit of etale neighbourhoods as the colimit of the cohomologies. Supplies the identification identifying the colimit of the sections (respectively the cohomology groups) over the approximating etale schemes with the corresponding group over the limit.
\item Covering of an object of the pro-etale site by limits of affine etale schemes. Same chapter, the proposition on weakly etale morphisms providing such coverings, and the lemma computing finite limits of sections. Supply the reduction to the case where the object is a limit V = lim V\_i of affine schemes etale over Y, and the description of the left hand side of the map.
\end{itemize}\medskip
\ifdefined\hypertarget\hypertarget{sourcepointer-1}{}\fi
\noindent\textit{Source reference.} Original lemma 03Q6: sources/stacks-chow/tags.txt, source lines 4883--4883. Prerequisite pointer invoked at proof line(s) 4953; it names a standard fact of a sibling chapter that the counted argument uses, and every step of the counted proof built on it is retained verbatim in the delivered body.\par\medskip
\ifdefined\hypertarget\hypertarget{sourcepointer-2}{}\fi
\noindent\textit{Source reference.} Original reference 0GLZ: sources/stacks-proetale/pinned-original-source.tex, source lines 3977--3977. Prerequisite pointer invoked at proof line(s) 4946; it names a standard fact of the same chapter that the counted argument uses, and every step of the counted proof built on it is retained verbatim in the delivered body.\par\medskip
\ifdefined\hypertarget\hypertarget{sourcepointer-3}{}\fi
\noindent\textit{Source reference.} Original reference 099S: sources/stacks-proetale/pinned-original-source.tex, source lines 4058--4058. Prerequisite pointer invoked at proof line(s) 4945; it names a standard fact of the same chapter that the counted argument uses, and every step of the counted proof built on it is retained verbatim in the delivered body.\par\medskip
\ifdefined\hypertarget\hypertarget{sourcepointer-4}{}\fi
\noindent\textit{Source reference.} Original reference 097Z: sources/stacks-proetale/pinned-original-source.tex, source lines 1802--1802. Prerequisite pointer invoked at proof line(s) 4935; it names a standard fact of the same chapter that the counted argument uses, and every step of the counted proof built on it is retained verbatim in the delivered body.\par\medskip
\ifdefined\hypertarget\hypertarget{sourcepointer-5}{}\fi
\noindent\textit{Source reference.} Original section 06UM: sources/stacks-chow/tags.txt, source lines 8751--8751. Prerequisite pointer invoked at proof line(s) 4930; it names a standard fact of a sibling chapter that the counted argument uses, and every step of the counted proof built on it is retained verbatim in the delivered body.\par\medskip
\begin{lemma}
\label{lemma-morphism-comparison}
Let $f : X \to Y$ be a morphism of schemes which is quasi-compact and
quasi-separated.
\begin{enumerate}
\item Let $\mathcal{F}$ be a sheaf of sets on $X_\etale$. Then we have
$f_{\proetale, *}\epsilon^{-1}\mathcal{F} =
\epsilon^{-1}f_{\etale, *}\mathcal{F}$.
\item Let $\mathcal{F}$ be an abelian sheaf on $X_\etale$. Then we have
$Rf_{\proetale, *}\epsilon^{-1}\mathcal{F} =
\epsilon^{-1}Rf_{\etale, *}\mathcal{F}$.
\end{enumerate}
\end{lemma}

\begin{proof}
Proof of (1). Let $\mathcal{F}$ be a sheaf of sets on $X_\etale$. There
is a canonical map $\epsilon^{-1}f_{\etale, *}\mathcal{F} \to
f_{\proetale, *}\epsilon^{-1}\mathcal{F}$, see
Sites, Section \ref{sites-section-pullback}.
To show it is an isomorphism we may work (Zariski) locally on $Y$, hence
we may assume $Y$ is affine. In this case
every object of $Y_\proetale$ has a covering by objects $V = \lim V_i$
which are limits of affine schemes $V_i$ \'etale over $Y$ (by
Proposition \ref{proposition-weakly-etale}
for example). Evaluating the map
$\epsilon^{-1}f_{\etale, *}\mathcal{F} \to
f_{\proetale, *}\epsilon^{-1}\mathcal{F}$
on $V$ we obtain a map
$$
\colim \Gamma(X \times_Y V_i, \mathcal{F})
\longrightarrow
\Gamma(X \times_Y V, \epsilon^*\mathcal{F})
$$
See Lemma \ref{lemma-limit-pullback} for the left hand side.
By Lemma \ref{lemma-describe-pullback-from-etale} we have
$$
\Gamma(X \times_Y V, \epsilon^*\mathcal{F}) =
\Gamma(X \times_Y V, g_\etale^{-1}\mathcal{F})
$$
where $g : X \times_Y V \to X$ is the projection. Hence the result holds by
\'Etale Cohomology, Lemma
\ref{etale-cohomology-lemma-directed-colimit-cohomology}.

\medskip\noindent
Proof of (2). Arguing in exactly the same manner as above
we see that it suffices to show that
$$
\colim H^i_\etale(X \times_Y V_i, \mathcal{F})
\longrightarrow
H^i_\etale(X \times_Y V, \mathcal{F})
$$
which follows once more from \'Etale Cohomology, Lemma
\ref{etale-cohomology-lemma-directed-colimit-cohomology}.
\end{proof}

\end{document}
